\section{What the posterior costs}
\label{sec:cost}

% Four results, in this order: the posterior's cost is fixed per step, its share
% therefore falls (with the small-B caveat in the main text, not a footnote), the
% part of it that is redundant traffic is removable, and the batch size decides
% which of the last two a reader is in. Interpret the tables; do not retype them.

\subsection{The cost is fixed per step}
\label{sec:cost:flat}

Each arm's posterior work costs the same milliseconds at $64$ tokens per step as at
$16{,}384$: across the $256$-fold sweep the directly timed Bayesian regions hold to
within $1.5\%$ of themselves on every arm (Table~\ref{tab:posterior}), the width of
the measurement and not a trend in it.

Nothing in those regions has a token in it: they stream parameter-sized buffers, and
how many tokens produced a gradient does not change its length. The traffic
is what the time is, since every parameter-sized region of the mean-field step and
their aggregate land in one bandwidth band (Table~\ref{tab:bandwidth}), so regions
sharing no arithmetic differ only in how many buffers each walks.

\begin{table}[t]
  \centering
  \resizebox{\linewidth}{!}{\input{\generated/tab-posterior}}
  \caption{Each arm's own Bayesian work, timed at the launch sites, unfused. Share is
  the posterior over the whole step, $k$ the exponent in $\text{share} \propto B^{-k}$
  between the two smallest and the two largest points measured.}
  \label{tab:posterior}
\end{table}

\subsection{Amortisation, and where it stops}
\label{sec:cost:amort}

The network's cost grows with tokens and the posterior's does not, so the
posterior's share of a step falls, at a local exponent of $1.005$ for mean-field
over the last two points (Table~\ref{tab:posterior}).

That asymptote is not the whole curve. Between $B = 1$ and $B = 2$ the exponent is
$0.148$ for mean-field and $0.196$ for Vadam, and the reason is the denominator: the
network is not $O(B \cdot T)$ at the bottom of the sweep either. The control arm
shows it with no posterior in the way: twice the tokens costs it a factor of $1.16$
there and a factor of two from $B = 16$ up (Table~\ref{tab:amortization}). Below that
a forward and backward pass costs more in launches than in tokens, so $1/B$ is the
asymptote and not the curve.

The fixed numerator is not a property of posteriors either: the control's own
parameter-sized share falls just as far over the same sweep
(Table~\ref{tab:amortization}). What fails to amortise is state not proportional to
the parameter count, which is why the low-rank arm still owns more than a tenth of a
step at its largest measurable batch. All of it is per-step cost, pricing neither the
ensemble's $E$ runs nor a variational arm's $S$ evaluation passes.

\begin{table}[t]
  \centering
  \resizebox{\linewidth}{!}{\input{\generated/tab-amortization}}
  \caption{The batch sweep for the shipped code. $O(P)\%$ is all parameter-sized work
  over the step, Bayesian or not. The ensemble column is one member, no posterior.}
  \label{tab:amortization}
\end{table}

\subsection{The removable part is the redundant traffic}
\label{sec:cost:fusion}

Folding an arm's parameter-sized passes into one kernel buys time in proportion to the
traffic it deletes, and nothing where there is none to delete. Mean-field's $32$
touches drop to $15$ for $1.76\times$ of $O(P)$ time at $B = 1$, the low-rank arm's
$122$ per layer to $58$ for $1.14$ to $1.16\times$, and Vadam's $11$ to $9$ for
between $0.98$ and $1.04\times$
(Table~\ref{tab:fusion}). The control is unchanged code in both sweeps and comes out
within $1\%$ of itself, so anything inside that band is drift, which is where
Vadam's sits.

That null is what Vadam's pass count predicts. Two touches of slack against a floor
of $9$ is $1.22\times$ of traffic, and a fused kernel gives back bandwidth:
mean-field's streams eight buffers at once where the sequence it replaced walked three
or four each (Table~\ref{tab:bandwidth}). At that discount a $2.13\times$ traffic cut
is still $1.76\times$ of time and a $1.22\times$ cut is nothing.

\begin{table}[t]
  \centering
  \resizebox{\linewidth}{!}{\input{\generated/tab-fusion}}
  \caption{Parameter-sized work before and after folding, where both sweeps measured
  the point; the last column is the share after over the share before. The ensemble's
  code is identical in both, so its rows are the drift.}
  \label{tab:fusion}
\end{table}

\subsection{Which of the two applies}
\label{sec:cost:crossover}

% The rule the last two subsections imply together, and the batch size where it
% turns over. It closes them against each other: a share that amortises and a fold
% worth whole multiples are one measurement, not two claims.

Folding deletes a fixed count of buffer touches, so what it returns is a fixed
wall-clock, and a fixed quantity amortises exactly as the posterior's own cost does.
The ratio is scale-free and the milliseconds are not: the same $1.76\times$ is a fifth
of a mean-field step at $B = 1$ and under a percent of one at $B = 128$
(Table~\ref{tab:crossover}). A share that amortises to nothing and a fold worth whole
multiples are one measurement at two batch sizes.

Taking the boundary per arm at the $1\%$ the control reproduces itself in, below which
a saving is not separable from the card's temperature between sweeps,
mean-field crosses at $B = 128$, Vadam is inside that band at every point measured, and
the low-rank arm is outside it wherever this card can run, so it has no regime where
folding fails to pay. Which side a run sits on is how large a step it takes, and a step
that fills this card is above mean-field's crossover: three arms reach $16{,}384$ tokens
per step on $6$\,GB, twice it. Below lie a batch small enough to be launch-bound and a
posterior whose own footprint holds the batch there.

\begin{table}[t]
  \centering
  \resizebox{\linewidth}{!}{\input{\generated/tab-crossover}}
  \caption{Where folding stops being separable from drift. What it returns is $O(P)$
  time deleted, over the unfused step; $B^\ast$ is the first batch size measured at
  which that lands inside the control's band, or a bound where the sweep never reaches
  or never leaves it.}
  \label{tab:crossover}
\end{table}
